\documentclass[11pt]{article}
\usepackage[margin=1in]{geometry}
\usepackage{amsmath,amssymb,amsthm}
\usepackage[T1]{fontenc}
\usepackage{lmodern}
\usepackage{hyperref}
\hypersetup{colorlinks=true,linkcolor=blue,urlcolor=blue}
\newtheorem{theorem}{Theorem}
\newtheorem{lemma}{Lemma}
\newcommand{\R}{\mathbb R}
\newcommand{\N}{\mathbb N}
\title{Disproof of Conjecture 00000000214\\\large Positive third differences of the partition logarithm occur arbitrarily late}
\author{Personal submission}
\date{}
\begin{document}
\maketitle
\begin{abstract}
Let $p(n)$ be the number of integer partitions of $n$ and put $f(n)=\log p(n)$.
For every $N$ there is an $n\ge N$ such that $\Delta^3 f(n)>0$.
Consequently, the third difference is not eventually negative.
The proof uses a finite encoding of actual partitions and an elementary principle for real sequences. It uses neither a partition asymptotic formula nor a finite numerical test of an eventual assertion. All mathematical steps are formalized in Lean 4.19.0 with Mathlib v4.19.0.
\end{abstract}

\section{Statement and scope}
The supplied English statement reads:
\begin{quote}
Definition: The third difference of log p(n). Conjecture: The third difference is eventually negative, a strengthening of the proven log-concavity; and the negativity is settled by the convexification layer. (partition log third difference negative)
\end{quote}
The exact bilingual source is included as \texttt{source/ORIGINAL.md}; its SHA-256 is
\begin{center}\small\texttt{101ccf808d4022fc597cef8e3c222960742ea86822643ef052f6b1cb62437347}.\end{center}
Here an integer partition is a finite multiset of positive integers whose sum is $n$. We use the standard forward difference
\[
\Delta f(n)=f(n+1)-f(n),\qquad
\Delta^3f(n)=f(n+3)-3f(n+2)+3f(n+1)-f(n).
\]
The source's necessary eventual-negativity assertion is
\begin{equation}\label{eq:claim}
\exists N\in\N\;\forall n\in\N\quad
 n\ge N\ \Longrightarrow\ \Delta^3\log p(n)<0.
\end{equation}
We disprove this assertion for the actual partition-count sequence. The standard backward third difference at $n\ge3$ equals the forward third difference at $n-3$, so its eventual-negativity assertion is disproved as well.
The phrase ``convexification layer'' has no definition in the supplied source. No interpretation of that phrase is introduced or used: the failure of the necessary assertion \eqref{eq:claim} suffices to refute the stated conjecture. Nor is any log-concavity result assumed.

\begin{theorem}\label{thm:main}
For every $N\in\N$, there exists $n\ge N$ for which
\[
\log p(n+3)-3\log p(n+2)+3\log p(n+1)-\log p(n)>0.
\]
In particular, \eqref{eq:claim} is false.
\end{theorem}

\section{Elementary bounds for actual partitions}
There is at least one partition of every $n$, including the empty partition of zero. Thus $p(n)\ge1$ and $f(n)\ge0$.
For $n>0$ and each $i\in\{0,\ldots,n-1\}$, consider the partition consisting of one part $i+1$ and $n-i-1$ parts equal to $1$. Its number of parts is $n-i$, so these $n$ partitions are distinct. Therefore
\begin{equation}\label{eq:lower}
p(n)\ge n\qquad(n\in\N).
\end{equation}
It follows that $f$ is unbounded above: for any real $C$, choose $n>e^C$ and use $p(n)\ge n$.

\begin{lemma}\label{lem:count}
If $n,k\in\N$ and $n\le k^2$, then
\[
p(n)\le(n+1)^{k+1}(k+1)(n+1)^k.
\]
\end{lemma}
\begin{proof}
Every part of a partition of $n$ is at most $n$, and its total number of parts is at most $n$. Call a part small when it is at most $k$, and large otherwise.
Every large part is at least $k+1$; hence there can be at most $k$ large parts, since $k+1$ of them would have sum at least $(k+1)^2>n$.
Encode a partition by:
\begin{enumerate}
\item the multiplicities of each of $0,1,\ldots,k$, each in $\{0,\ldots,n\}$;
\item the number of large parts, in $\{0,\ldots,k\}$;
\item the large parts in increasing order, padded with zeros to a list of length $k$, each entry in $\{0,\ldots,n\}$.
\end{enumerate}
This encoding is injective. The stored number of large parts identifies the length of the unpadded sorted list, and its entries recover the large-part multiset. The small-part multiplicities recover the other multiset. Their union is the original partition. The codomain has the displayed cardinality. This also covers $k=0$, which forces $n=0$.
\end{proof}
Taking logarithms of Lemma~\ref{lem:count}, with positive arguments, gives
\begin{equation}\label{eq:logbound}
f(n)\le(2k+1)\log(n+1)+\log(k+1)\qquad(n\le k^2).
\end{equation}
As $k^2+1\le(k+1)^2$, in particular
\begin{equation}\label{eq:square}
f(k^2)\le(4k+3)\log(k+1).
\end{equation}

\begin{lemma}\label{lem:noaffine}
For every $a>0$ and every $b\in\R$, there exists $n\in\N$ such that
\[
f(n)<b+an.
\]
\end{lemma}
\begin{proof}
The elementary real limit $\log x/x\to0$ as $x\to\infty$ implies that there is an $M\in\R$ such that
\[
|\log x|\le\frac a{28}|x|\qquad(x\ge M).
\]
Choose an integer $k>\max\{M,1,-2b/a\}$. Then
$\log(k+1)\le a(k+1)/28$, and
$(4k+3)(k+1)\le14k^2$. Equation~\eqref{eq:square} yields
\[
f(k^2)\le\frac a2 k^2.
\]
Since $-2b<ak\le ak^2$, this is strictly smaller than $b+ak^2$. Take $n=k^2$.
The cited elementary limit is the stock Mathlib theorem
\texttt{Real.isLittleO\_log\_id\_atTop}, applied through its explicit eventual bound. No asymptotic fact about $p$ is imported.
\end{proof}

\section{The sequence obstruction}
\begin{lemma}\label{lem:sequence}
Let $f:\N\to\R$ be nonnegative and unbounded above. Suppose that for every $a>0$ and every $b\in\R$ there exists $n$ such that $f(n)<b+an$. Then its third differences cannot be nonpositive throughout a tail.
\end{lemma}
\begin{proof}
We first record two elementary facts about a nonnegative sequence $u$.
If $\Delta u(j)\le A$ for all $j$, telescoping gives
$u(j)\le u(0)+jA$. This forces $A\ge0$, since otherwise the right side is negative for some integer $j$.
Consequently, if $\Delta u$ is nonincreasing, then $\Delta u(m)\ge0$ for each $m$: apply the preceding fact to the tail $u(m+j)$ with upper bound $A=\Delta u(m)$.

Now assume $\Delta^2u$ is nonincreasing. It must be nonnegative everywhere. Indeed, if $\Delta^2u(m)<0$, then on the tail starting at $m$ the first differences are nonincreasing, so they are nonnegative by the preceding paragraph. But this nonnegative first-difference sequence has all its first differences bounded above by the negative constant $\Delta^2u(m)$, a contradiction to the same elementary fact.

Suppose for contradiction that $\Delta^3f(n)\le0$ for all $n\ge N$, and let $u(j)=f(N+j)$. Then $\Delta^2u$ is nonincreasing. By the previous paragraph, $\Delta^2u\ge0$, so $\Delta u$ is nondecreasing.
If $\Delta u(k)=a>0$ for some $k$, write $M=N+k$. Every first difference on the tail beginning at $M$ is at least $a$. Telescoping and nonnegativity imply
\[
f(M+j)\ge ja\qquad(j\in\N).
\]
For $n\le M$, nonnegativity gives $f(n)\ge0\ge a(n-M)$, and for $n\ge M$ the displayed inequality gives the same lower bound. Hence
$f(n)\ge-aM+an$ for every $n$, contrary to the hypothesis with $b=-aM$.
Thus $\Delta u\le0$ everywhere. The tail $u$ is nonincreasing and bounded above by $u(0)$. A finite prefix is bounded as well, so $f$ is bounded above, contradicting the other hypothesis.
\end{proof}
Applying Lemma~\ref{lem:sequence} to the actual partition logarithm, using \eqref{eq:lower} and Lemma~\ref{lem:noaffine}, proves
\[
\neg\exists N\;\forall n\ge N,\quad \Delta^3\log p(n)\le0.
\]
Classical negation yields Theorem~\ref{thm:main}. This proves arbitrarily late strict positivity; it makes no claim that positivity is eventual and gives no numerical cutoff.

\section{Formal correspondence and verification}
The formal definition is
\[
\begin{aligned}
\texttt{partitionCount n}&=\texttt{Fintype.card (Nat.Partition n)},\\
\texttt{partitionLog n}&=\texttt{Real.log (partitionCount n :}\;\R\texttt{)}.
\end{aligned}
\]
Mathlib's \texttt{Nat.Partition n} is precisely a multiset of positive natural numbers with sum $n$. Thus the theorem binds the actual mathematical objects, with no growth, monotonicity, or recurrence assumption left as a hypothesis.

\begin{center}
\begin{tabular}{ll}
\textbf{File} & \textbf{Role}\\\hline
\texttt{TLMC214/PartitionBounds.lean} & Actual partitions, injections, cardinality bounds\\
\texttt{TLMC214/Growth.lean} & Logarithms, unboundedness, Lemma~\ref{lem:noaffine}\\
\texttt{TLMC214/Sequence.lean} & Generic sequence Lemma~\ref{lem:sequence}\\
\texttt{TLMC214/Disproof.lean} & Concrete negation and strict-positive witnesses
\end{tabular}
\end{center}
The theorem \texttt{TLMC214.disproof} has conclusion
\texttt{\(\neg\) TLMC214.EventualNegativity}.
The theorem \texttt{TLMC214.positive\_arbitrarily\_late} quantifies over every cutoff $N$ and supplies a later index with strictly positive third difference. The theorem \texttt{TLMC214.backward\_disproof} handles the standard backward convention by an index shift.

The project pins Lean 4.19.0 and Mathlib commit
\begin{center}\small\texttt{c44e0c8ee63ca166450922a373c7409c5d26b00b}.\end{center}
The build treats warnings as errors. Separate source replays and the \texttt{CheckAxioms.lean} audit are described in the README and retained in the evidence directory. The capstone declarations depend only on the usual Lean foundations \texttt{propext}, \texttt{Classical.choice}, and \texttt{Quot.sound}. There are no admitted goals, custom axioms, native-decision shortcuts, numerical experiments, or auxiliary numerical computations.

The exact input, fresh mathematical development, failed proof attempts, command receipts, and final source hashes are retained for provenance. Neither maintainer acceptance nor publication eligibility is asserted by this report.
\end{document}
